% Source: https://artofproblemsolving.com/community/c3h2916240_2022_mbmt_guts_round_d115_z18_montgomery_blair_math_tournament
% Stable identity: aops-topic-2916240-q13
Ten consecutive integers contain at least eight factors of 2, so the number of zeros is controlled by factors of 5 here. If the interval ends before 25, its two multiples of 5 each contribute one factor of 5, yielding only two zeros. The first interval reaching 25 begins at 16. Its multiples 20 and 25 contribute $1+2=3$ factors of 5, so it has exactly three zeros. Therefore $\boxed{16}$.
